\section{大模型对计算系统提出的要求}
\label{sec:req}

\begin{guide}
\guideitem{需求}{在讨论任何硬件或软件之前，首先明确需要计算多少、存储多少、搬运多少。}
\guideitem{方法}{用参数量 $N$、每个 token 的计算量和每次计算的数据搬运量做数量级估算，再用 Roofline 模型判断瓶颈。}
\guideitem{结论}{大模型同时受到算力、存储容量、存储带宽和设备互联四方面约束，由此形成计算系统的四层结构。}
\end{guide}

\subsection{两个典型场景}

先看两个场景。

\emph{场景一：训练。}某团队希望用 $2\times10^{12}$ 个 token 的语料，从头训练一个 70 亿参数（7B）的模型。他们需要知道：单张显卡的显存能否容纳模型？训练需要多长时间？需要多少张显卡？

\emph{场景二：推理。}训练好的模型部署为在线问答服务，同时有上百名用户提问。服务方关心的是：用户需要等待多久才能看到第一个输出 token？每秒能输出多少个 token？一台服务器能同时服务多少用户？

这两个问题看似不同，背后却涉及同样的三个基本量：\emph{存储量}（模型与中间状态占多少字节）、\emph{计算量}（需要多少次浮点运算）和\emph{数据搬运量}（运算时要在存储与计算单元之间移动多少字节）。下面逐一估算。

\subsection{三个基本量的估算}
\label{subsec:three-quantities}

\paragraph{存储量} 设模型有 $N$ 个参数，每个参数用 $b$ 位表示，则仅权重就至少占用
\begin{equation}
  M_{\mathrm{weight}} = \frac{N\,b}{8}\ \text{字节}.
  \label{eq:weight-mem}
\end{equation}
7B 模型若以 16 位浮点数（FP16 或 BF16）存储，需要 $7\times10^{9}\times 2 = 14\times10^{9}$ 字节，约 \SI{14}{GB}；若量化为 4 位，仅权重约 \SI{3.5}{GB}。训练时还要额外保存梯度和优化器状态，所需存储往往是权重本身的数倍（详见第\,\ref{subsec:train-mem}\,小节）。

\paragraph{计算量} 对于 Transformer 结构的模型，处理一个 token 的前向计算约需 $2N$ 次浮点运算（每个参数参与一次乘法和一次加法）；训练时还要做反向传播，计算量约为前向的两倍，因此每个训练 token 合计约 $6N$ 次浮点运算\cite{kaplan2020scaling}。于是训练总计算量可估算为
\begin{equation}
  C_{\mathrm{train}} \approx 6\,N\,D,
  \label{eq:train-flops}
\end{equation}
其中 $D$ 为训练 token 总数。

\begin{example}[训练需要多少张卡]{ex:gpus}
按场景一，$C_{\mathrm{train}}\approx 6\times7\times10^{9}\times2\times10^{12}=8.4\times10^{22}$ 次浮点运算。一张数据中心级 GPU 的 BF16 稠密峰值算力约为 $10^{15}$\,FLOP/s 量级（例如 NVIDIA H100 SXM 约 $9.9\times10^{14}$\,FLOP/s\cite{nvidia2023h100}）。实际训练中能达到峰值的 40\% 左右已属较高水平，即有效算力约 $4\times10^{14}$\,FLOP/s，则
\[
  t_{\text{单卡}} \approx \frac{8.4\times10^{22}}{4\times10^{14}} \approx 2.1\times10^{8}\ \text{s} \approx 6.7\ \text{年}.
\]
若使用 1024 张 GPU 并行训练，且并行效率理想，耗时可缩短到约 2.4 天。这说明：\emph{大模型训练通常需要多卡乃至多机的分布式计算}，设备之间如何交换数据因此成为核心问题。
\end{example}

\paragraph{数据搬运量} 计算单元只能处理已经送达的数据。推理时每生成一个新 token，都要把全部权重从显存读入计算单元一次：对于 \SI{14}{GB} 权重的 7B 模型，就意味着每个 token 至少读取 \SI{14}{GB} 数据。若显存带宽为 \SI{3.35}{TB/s}，则单个请求逐 token 生成的速度上限约为
\[
  \frac{3.35\times10^{12}\ \text{B/s}}{14\times10^{9}\ \text{B/token}} \approx 240\ \text{token/s}.
\]
该上限主要不由芯片峰值算力决定，而由\emph{存储带宽}决定。可见，衡量计算系统时不能只看算力，还必须同时考察数据供给是否及时。

\subsection{瓶颈判断工具：算术强度与 Roofline 模型}
\label{subsec:roofline}

为判断一项计算究竟受算力限制还是受访存限制，引入\term{算术强度}（\en{arithmetic intensity}）：记一次计算的浮点运算次数为 $F$，需要在存储与计算单元之间搬运的字节数为 $Q$，则
\begin{equation}
  I = \frac{F}{Q}\quad(\text{FLOP/byte}).
\end{equation}
设硬件峰值算力为 $P_{\mathrm{peak}}$，存储带宽为 $B_{\mathrm{mem}}$。计算单元每秒最多完成 $P_{\mathrm{peak}}$ 次运算，而存储系统每秒最多供给 $B_{\mathrm{mem}}$ 字节数据、支撑 $I\cdot B_{\mathrm{mem}}$ 次运算，二者取较小者，即得 \term{Roofline 模型}\cite{williams2009roofline}给出的可达性能上界：
\begin{equation}
  P_{\mathrm{attainable}} \le \min\bigl(P_{\mathrm{peak}},\; I\cdot B_{\mathrm{mem}}\bigr).
  \label{eq:roofline}
\end{equation}
两条直线相交处的算术强度 $I^{*}=P_{\mathrm{peak}}/B_{\mathrm{mem}}$ 称为\term{拐点}。$I<I^{*}$ 时性能受带宽限制，称为\term{访存受限}（\en{memory-bound}）；$I>I^{*}$ 时性能受峰值算力限制，称为\term{计算受限}（\en{compute-bound}）。以前述 GPU 为例，$I^{*}\approx 9.9\times10^{14}/3.35\times10^{12}\approx 295$\,FLOP/byte。

\begin{figure}[htbp]
\centering
\begin{tikzpicture}[fig]
\begin{axis}[
  width=0.86\linewidth, height=6.2cm,
  xmode=log, ymode=log, log basis x=10, log basis y=10,
  xmin=0.3, xmax=10000, ymin=1, ymax=3000,
  xlabel={算术强度 I / (FLOP/byte)},
  ylabel={可达性能 / TFLOP/s},
  label style={font=\figfont\small}, tick label style={font=\figfont\small},
  axis line style={figline}, tick style={figline},
  xtick={1,10,100,1000,10000}, ytick={1,10,100,1000},
  xticklabels={1,10,100,{1\,000},{10\,000}}, yticklabels={1,10,100,{1\,000}},
  minor tick num=0, clip=false,
]
  % 屋顶线：min(989, 3.35 I)
  \addplot[draw=ink, line width=1.1pt, domain=0.3:295, samples=2] {3.35*x};
  \addplot[draw=ink, line width=1.1pt, domain=295:10000, samples=2] {989};
  % 拐点
  \draw[dashed, rulegray] (axis cs:295,1) -- (axis cs:295,989);
  \node[note, anchor=south west, inner sep=2pt] at (axis cs:360,1.5) {拐点 I\textsuperscript{*} ≈ 295};
  \node[note, anchor=south] at (axis cs:2300,1050) {计算受限区：P\textsubscript{peak}};
  \node[note, anchor=south west] at (axis cs:0.4,40) {访存受限区：I\,·\,B\textsubscript{mem}};
  % 解码（批大小 1）
  \addplot[only marks, mark=*, mark size=2.4pt, warn] coordinates {(1,3.35)};
  \node[note, anchor=west, text=warn] at (axis cs:1.25,2.4) {解码，B = 1：I ≈ 1};
  % 解码（批大小 64）
  \addplot[only marks, mark=*, mark size=2.4pt, figline] coordinates {(64,214)};
  \node[note, anchor=east] at (axis cs:52,240) {解码，B = 64：I ≈ 64};
  % 预填充
  \addplot[only marks, mark=*, mark size=2.4pt, accent] coordinates {(2048,989)};
  \node[note, anchor=north, text=accent] at (axis cs:2048,700) {预填充，T = 2048};
\end{axis}
\end{tikzpicture}
\caption{Roofline 模型示意。以峰值算力约 \SI{989}{TFLOP/s}、显存带宽约 \SI{3.35}{TB/s} 的 GPU 为例；坐标为对数刻度，各点为数量级估计}
\label{fig:roofline}
\end{figure}


将 Roofline 模型用于推理的两个阶段（图\,\ref{fig:roofline}），可以得到一个重要结论。读入一次 16 位权重（2 字节）后，若只为一个 token 做 2 次运算，算术强度约为 $1$\,FLOP/byte，远低于拐点；若一次读入后同时为 $T$ 个 token 复用，算术强度约为 $T$。因此：
\begin{itemize}
  \item \emph{预填充}（一次性处理用户输入的全部 token，$T$ 可达数千）的算术强度高，通常\emph{计算受限}；
  \item \emph{逐 token 解码}（每步只生成一个 token）的算术强度很低，通常\emph{访存受限}；把多个请求合并成批次（批大小为 $B$）可使算术强度提高到约 $B$。
\end{itemize}
这就解释了一个常见现象：同一张显卡处理长输入（预填充阶段）时吞吐较高，而逐 token 解码时吞吐却较低——原因不是算力不足，而是每生成一个 token 都要重新读取大量权重和缓存。第\,\ref{sec:infer}\,节介绍的许多推理优化技术，正是围绕这一点展开的。

\begin{pitfall}
Roofline 模型给出的是性能\emph{上界}，不能用来精确预测实际性能。它的价值在于快速判断主要瓶颈：对访存受限的任务，提升峰值算力作用有限；对计算受限的任务，提升带宽同样作用有限。
\end{pitfall}

\subsection{从需求到方案：计算系统的四层结构}
\label{subsec:four-layers}

综合以上估算，大模型对计算系统的需求可以归纳为四条：
\begin{enumerate}
  \item \emph{算得快}——提供足够的浮点算力，并让算力被充分利用；
  \item \emph{存得下、送得及}——有足够的存储容量容纳权重和中间状态，并有足够的带宽及时供给数据；
  \item \emph{传得动}——当一张卡不够用时，多张卡、多台机器之间能高效交换数据；
  \item \emph{用得上}——开发者能方便地表达模型、训练模型，用户能方便地调用模型。
\end{enumerate}

没有哪一个部件能单独满足这四条需求，于是形成了层次化的解决方案（图\,\ref{fig:overview}）：底层的\emph{算力资源}（芯片、存储、互联）提供物理能力；\emph{驱动与计算软件栈}把硬件能力开放为可编程接口；\emph{训练与推理框架}让开发者用张量运算表达模型，并自动完成求导、并行与调度；最上层的\emph{应用系统工具}把模型组装成可部署、可管理的产品。每一层都向上层隐藏下层的复杂性，同时受到下层能力的约束。

\begin{figure}[htbp]
\centering
\begin{tikzpicture}[fig, x=1mm, y=1mm]
  \def\W{92}
  % 层
  \node[blk, minimum width=\W mm, minimum height=13mm] (L1) at (0,0)
    {{\bfseries 应用系统工具}\\[1pt]\footnotesize 编排平台、应用开发框架、向量数据库、微调工具、部署工具};
  \node[blk, minimum width=\W mm, minimum height=13mm] (L2) at (0,-21)
    {{\bfseries 训练与推理框架}\\[1pt]\footnotesize PyTorch、DeepSpeed / Megatron-LM、vLLM / SGLang};
  \node[blk, minimum width=\W mm, minimum height=13mm] (L3) at (0,-42)
    {{\bfseries 驱动与计算软件栈}\\[1pt]\footnotesize 编译器、运行时、算子库、通信库（CUDA / CANN / NeuWare）};
  % 硬件层：三个子块
  \node[grp, minimum width=\W mm, minimum height=24mm] (L4) at (0,-68) {};
  \node[anchor=north, font=\figfont\small\bfseries] at (L4.north) {算力资源};
  \node[blkw, minimum width=27.5mm, minimum height=12mm] (C) at (-30.5,-71) {计算芯片\\\footnotesize CPU、GPU、NPU};
  \node[blkw, minimum width=27.5mm, minimum height=12mm] (M) at (0,-71) {存储\\\footnotesize HBM、DRAM、SSD};
  \node[blkw, minimum width=27.5mm, minimum height=12mm] (N) at (30.5,-71) {互联\\\footnotesize PCIe、NVLink、IB};
  % 层间箭头与说明
  \foreach \a/\b/\t in {L1/L2/调用模型与服务, L2/L3/张量算子与计算图, L3/L4/内核执行与数据搬运} {
    \draw[flow] ([xshift=-30mm]\a.south) -- ([xshift=-30mm]\b.north)
      node[midway, right, note] {\t};
  }
  \draw[flow, <-] ([xshift=30mm]L1.south) -- ([xshift=30mm]L2.north) node[midway, left, note] {提供训练与推理能力};
  \draw[flow, <-] ([xshift=30mm]L2.south) -- ([xshift=30mm]L3.north) node[midway, left, note] {提供可编程接口};
  \draw[flow, <-] ([xshift=30mm]L3.south) -- ([xshift=30mm]L4.north) node[midway, left, note] {提供算力、容量与带宽};
  % 需求标签（左侧）
  \foreach \n/\t in {L1/用得上, L2/易表达、可扩展, L3/发挥硬件效率} {
    \node[note, anchor=east, align=right] at ([xshift=-3mm]\n.west) {\t};
  }
  \node[note, anchor=east, align=right] at ([xshift=-3mm]L4.west) {算得快\\存得下\\传得动};
  \node[note, anchor=south, font=\figfont\footnotesize\bfseries] at (-58,8) {需求};
  % 章节对应（右侧）
  \foreach \n/\t in {L1/4.5\,节, L2/4.3、4.4\,节, L3/4.2.4\,小节, L4/4.2.1\textendash{}4.2.3\,小节} {
    \node[note, anchor=west, align=left] at ([xshift=3mm]\n.east) {\t};
  }
  \node[note, anchor=south, font=\figfont\footnotesize\bfseries] at (58,8) {本章位置};
\end{tikzpicture}
\caption{大模型计算系统的四层结构。左侧为各层满足的需求，右侧为对应的章节；向下的箭头表示上层向下层提交的工作，向上的箭头表示下层向上层提供的能力}
\label{fig:overview}
\end{figure}


表\,\ref{tab:roadmap} 给出了本章各节与这四层结构的对应关系。阅读时可以随时回到图\,\ref{fig:overview}，确认当前内容处在哪一层、解决的是哪一类需求。

\begin{table}[htbp]
\centering
\caption{本章内容与系统层次的对应关系}
\label{tab:roadmap}
\small
\begin{tabularx}{\linewidth}{@{}l l Y@{}}
\toprule
\heiti 节 & \heiti 所处层次 & \heiti 回答的核心问题 \\
\midrule
\ref{sec:hw}\quad 算力资源 & 算力资源、驱动与软件栈 & 计算在哪里执行？数据如何存取和传输？ \\
\ref{sec:fw}\quad 深度学习框架 & 训练与推理框架 & 如何方便地表达模型？单卡放不下时如何训练？ \\
\ref{sec:infer}\quad 推理服务 & 训练与推理框架 & 如何让模型又快又省地回答大量请求？ \\
\ref{sec:tools}\quad 应用系统工具 & 应用系统工具 & 如何把模型变成可使用、可管理的产品？ \\
\ref{sec:case}\quad 综合案例 & 全部层次 & 面对一个具体需求，如何自上而下地完成选型？ \\
\bottomrule
\end{tabularx}
\end{table}
